% ===== §2 Thesis: generativity takes its own continuation as its end =====
\section{Thesis: Generativity Takes Its Own Continuation as Its End}
\label{p09:sec:thesis}

\noindent
The paper begins from a metaphysical conjecture, and states it at once, baldly, so that the antithesis may have something firm to push against. \emb{Generativity takes its own continuation as its end. Production points back at production itself.} The aim of a generative process is not some terminal product at which it would come to rest, but the perpetuation of generating---the cycle's end is the cycle. A generativity oriented toward a final product would, upon attaining it, stop; the generativity this paper is about does not aim at a state it could reach and rest in, but at its own ongoingness. Its telos is itself.

This is a strong claim, and its strangeness should not be smoothed over. We are accustomed to think of processes as means to ends---production as the making of a product, desire as the seeking of an object, labour as the securing of a wage---and on that picture the process is subordinate, justified by the end it serves, and complete when the end is attained. The conjecture inverts this: in the generativity that concerns us, the ``product'' is not the point at which the process is justified and stops, but the occasion for the process to continue. The product is for the sake of the producing, not the producing for the sake of the product. This is what is meant by calling such generativity \emph{autotelic}: it carries its end within itself, as its own continuation, and is for that reason not completed by any external attainment.

The conjecture is not arbitrary, and its warrant is the constancy with which it recurs across domains that share no subject matter. Four such recurrences are worth setting out, for the breadth of the pattern is the chief reason to take it for something like a metaphysical ground rather than a local fact about one or another region of the world.

\emc{In the dynamics of the cosmos}, the structures that persist are, overwhelmingly, the self-reproducing ones. A dissipative structure---a flame, a whirlpool, a living cell---maintains itself far from thermodynamic equilibrium not by reaching some final stable state but by continually regenerating the very gradients and flows that constitute it; it persists by producing the conditions of its own persistence. The autocatalytic cycle, the chemical paradigm of life's origin, is a set of reactions whose product is the catalyst of its own production---a loop whose output is the cause of its output. In both, what endures is precisely what makes more of its own enduring; the structures that aim at no terminal state, but only at their own continuation, are the ones the universe keeps. The autotelic is not an oddity of human striving but a feature of what persists at all.

\emc{In the dialectic of desire}, as psychoanalysis describes it, the same structure governs the subject. Desire does not aim at a final satisfaction that would extinguish it; were it ever fully satisfied, it would cease, and the subject of desire would dissolve. Beneath every particular object---this person, this attainment, this possession---desire aims at its own continuation, which is why the object attained is, with a regularity that is structural and not accidental, found wanting, and desire moves on. This is not the failure of desire but its success: desire is organized to defer the satisfaction that would end it, to keep itself going by never quite arriving. Lacan's distinction of desire from demand and need turns on just this: need has an object that satisfies and stops it, but desire, beneath demand, aims past every object at the perpetuation of desiring itself. The subject lives by a wanting whose true end is to go on wanting.

\emc{In the structure of dialectic itself}---thesis, antithesis, synthesis---the autotelic form appears as the engine of the movement. The synthesis that resolves a contradiction is not a resting place but the thesis of a further movement; the \emph{Aufhebung} that lifts a tension into a higher unity does not end the dialectic but generates its next turn. A dialectic that reached a final synthesis and stopped would not be a dialectic but a death; the living dialectic is the one whose every resolution opens the next contradiction, whose end is not a final synthesis but the continuation of synthesizing. Here, in the most abstract register, is the same structure the rose and the cell display: a process whose product is the occasion of its own continuation.

\emc{And in Eglash's generative justice}, this conjecture finds the political-economic and ethical articulation on which the paper will lean throughout. Eglash's insight is that value, in an unalienated economy, does not flow outward to be extracted and accumulated at a distant centre, but \emph{circulates back to those who generate it}, to be reinvested in the conditions of further generation \parencite{eglash2016}. The recursive, self-similar structure on which this regeneration leans is the through-line of Eglash's earlier study of indigenous fractal design \parencite{eglash1999}. Where the dominant economies---he argues that both the capitalist and the state-socialist are, in this, alike---are organized around the \emph{extraction} of value from its producers to a centre, the generative alternative keeps value circulating at the grassroots, among and back to those whose labour, care, and creativity produce it. Generativity, here, perpetuates itself by recursion: the value it produces returns to nourish the producers, who produce again, in a loop that is its own end. \emb{This is the spatial criterion of generative justice: that value, unalienated, return to its creators.} And the recursive structure of regeneration is, in Eglash's analysis, all but universal---generativity, to perpetuate itself at all, must lean on a recursive structure that returns its product to the conditions of its own renewal.

Drawing the four together:

\begin{claim}[The metaphysical ground, as point of departure]
Generativity takes its own continuation as its end---an autotelic, self-referential structure that recurs across the dynamics of the cosmos (the self-reproducing dissipative structure), the dialectic of desire (the wanting whose end is to go on wanting), the structure of dialectic itself (the synthesis that is the next thesis), and the political economy of generative justice (the recursive return of value to its creators). The constancy of the recurrence warrants treating it as something like a metaphysical ground. This is the paper's \emph{thesis}, its point of departure.
\end{claim}

But it must be stressed that the thesis is offered as ground and not as conclusion, for the temptation here is real and the whole argument of the paper turns on refusing it. To establish that generativity is self-perpetuating, and that justice consists in the recursive return of value to its creators, can seem already to have delivered an ethics: let value circulate back to those who make it, refuse the extraction that carries it off to a centre, and all is well. The labour of this paper is to show that this is not enough---that the recursive return of value is necessary to the good circle but does not suffice for it, because the bare fact of self-continuation, even self-continuation \emph{with} return, cannot by itself tell the good circle from the vicious one. For the most rapacious cycles are self-continuing too, and some of them return value to a circle of creators while exhausting everything outside it. To that insufficiency---the moral indifference of mere self-continuation---the antithesis now turns.
